\section{Backup and Recovery}

\begin{frame}{Why Backup and Recovery?}
  Data is one of the most valuable assets of any organisation. Hardware can be replaced in hours;
  data that is lost without a backup may be lost forever. Even a brief period of data unavailability
  can cause significant financial and reputational damage.

  \medskip
  Every production database system \textbf{must} have a well-defined backup and recovery strategy
  before going live -- not after the first incident.

  \medskip
  \begin{block}{Two key metrics for every recovery strategy}
    \begin{itemize}
      \item \textbf{RTO -- Recovery Time Objective}: the maximum acceptable duration of system
            downtime after a failure. ``How long can we afford to be offline?''
      \item \textbf{RPO -- Recovery Point Objective}: the maximum acceptable amount of data loss
            measured in time. ``How much data can we afford to lose?''
    \end{itemize}
  \end{block}

  \medskip
  Example: an online bank might have an RTO of 1 hour and an RPO of 0 minutes (no data loss
  permitted). A small internal tool might tolerate an RTO of 1 day and an RPO of 24 hours.
  Your backup strategy must be designed to meet these targets.
\end{frame}

\begin{frame}{Backups}
  \begin{block}{Definition}
    A \textbf{backup} is a copy of database data saved at a specific point in time. It is used to
    restore the database after data loss, corruption, accidental deletion, or hardware failure.
  \end{block}

  \medskip
  \textbf{Common backup strategies -- from coarsest to finest granularity:}

  \begin{itemize}
    \item \textbf{Full backup (e.g.\ weekly)}: a complete copy of the entire database. Reliable and
          self-contained -- restore from this alone is possible. Downside: slow to create and
          requires the most storage.
    \item \textbf{Differential backup (e.g.\ daily)}: copies only the data that has changed since
          the \emph{last full backup}. Faster and smaller than a full backup; restore requires the
          last full backup plus the most recent differential.
    \item \textbf{Incremental backup (e.g.\ hourly)}: copies only data changed since the
          \emph{last backup of any type}. Smallest and fastest; restore requires the full backup,
          every differential, and every incremental in sequence.
    \item \textbf{Transaction log backup}: saves a continuous log of every individual change. Enables
          \textbf{point-in-time recovery} -- restore to any moment, not just the last backup time.
  \end{itemize}

  \medskip
  \begin{examples}{Typical strategy in practice}
    Full backup on Sunday night + differential every other night + log backup every hour.
  \end{examples}
\end{frame}

\begin{frame}{Recovery}
  \begin{block}{Definition}
    \textbf{Recovery} is the process of restoring a database to a correct, consistent, and usable
    state after any type of failure.
  \end{block}

  \medskip
  \textbf{Common failure scenarios:}

  \begin{itemize}
    \item \textbf{Software issues}: application bugs that corrupt data, a transaction that completes
          partially before crashing, or a developer running \texttt{DELETE FROM orders} without a
          \texttt{WHERE} clause.
    \item \textbf{Hardware issues}: sudden power outage, disk head crash, server room fire or flood,
          failed storage controller.
    \item \textbf{Human error}: accidental \texttt{DROP TABLE}, updating the wrong rows, importing
          bad data that overwrites good data.
  \end{itemize}

  \medskip
  \textbf{Standard recovery process (3 steps):}
  \begin{enumerate}
    \item Restore the most recent \textbf{full backup} to get a consistent baseline.
    \item Apply the most recent \textbf{differential backup} to bring data closer to the failure
          point.
    \item Replay all \textbf{transaction log backups} in chronological order up to the desired
          recovery point in time.
  \end{enumerate}
\end{frame}

\begin{frame}{Transaction Logs: Before- and After-Images}
  To support both rollback and crash recovery, every database engine maintains detailed change logs
  for every transaction.

  \medskip
  \begin{block}{Before-Image}
    The state of a data record \textbf{before} the modification. Created at the start of a
    transaction (before any change is applied). Used for \texttt{ROLLBACK} -- if the transaction
    must be undone, the before-image is used to restore the original data.
    Kept in the log until the transaction either commits or rolls back.
  \end{block}

  \begin{block}{After-Image}
    The state of a data record \textbf{after} the modification. Created when the change is applied
    to the in-memory cache. Used for \textbf{REDO} / crash recovery -- if the system crashes after
    a commit but before the data was flushed to disk, the after-image allows the change to be
    replayed. Kept until the data has been physically written to disk (at a \emph{checkpoint}).
  \end{block}

  \medskip
  \textbf{Important:} data modifications are first written to a \textbf{database cache} (RAM buffer)
  for performance. They are written to physical disk asynchronously during a \emph{checkpoint}.
  If the system crashes between the commit and the next checkpoint, the transaction log (after-image)
  is used to redo the committed work.
\end{frame}

\begin{frame}{Transaction Lifecycle with Before- and After-Images}
  Concrete example: a transaction updates the price of ``Cat Milk'' from \EUR{1.99} to \EUR{2.49}.

  \medskip
  \begin{enumerate}
    \item \textbf{START TRANSACTION}: the DBMS begins tracking this transaction.
    \item \textbf{Create Before-Image}: the current record (price = 1.99) is written to the
          transaction log as the before-image.
    \item \textbf{Modify data in cache}: the price is updated to 2.49 in the in-memory buffer.
          The physical disk file is \emph{not yet changed}.
    \item \textbf{Create After-Image}: the updated record (price = 2.49) is written to the
          transaction log as the after-image.
    \item \textbf{COMMIT or ROLLBACK}:
      \begin{itemize}
        \item \textbf{COMMIT}: the change is confirmed. The before-image is deleted from the log.
              The after-image remains until the buffer is flushed to disk at the next checkpoint,
              then it too is deleted.
        \item \textbf{ROLLBACK}: the before-image is used to restore the record to price = 1.99.
              Both before- and after-image are discarded.
      \end{itemize}
  \end{enumerate}

  \medskip
  \begin{block}{Why this matters}
    This mechanism directly implements two ACID properties: \textbf{Atomicity} (all-or-nothing via
    rollback) and \textbf{Durability} (committed changes survive crashes via redo from the log).
  \end{block}
\end{frame}

\begin{frame}{RAID -- Redundant Array of Independent Disks (!)}
  \begin{block}{Definition (!)}
    \textbf{RAID} combines multiple physical hard drives into a single logical storage unit, providing
    a combination of improved performance (through parallelism) and/or fault tolerance (through
    redundancy), depending on the RAID level chosen.
  \end{block}

  \begin{itemize}
    \item Can be implemented in \textbf{hardware} (dedicated RAID controller card with its own CPU
          and cache) or in \textbf{software} (the operating system manages RAID in firmware or kernel
          drivers). Hardware RAID is faster and more reliable.
    \item \textbf{Three fundamental techniques} used in different combinations by each RAID level:
      \begin{itemize}
        \item \textbf{Striping}: data is split into blocks and written across multiple disks in
              parallel. Increases throughput -- both reading and writing happen simultaneously on
              several disks.
        \item \textbf{Mirroring}: an exact, real-time copy of all data is maintained on a second
              (or more) disk. If one disk fails, the mirror takes over immediately.
        \item \textbf{Parity}: a mathematical checksum of the data is stored alongside it. If one
              disk fails, its content can be \emph{reconstructed} from the remaining data and the
              parity information -- without needing a full mirror copy.
      \end{itemize}
    \item There is always a \textbf{trade-off} between performance, redundancy, and storage
          efficiency. No RAID level maximises all three simultaneously.
  \end{itemize}
\end{frame}

\begin{frame}{RAID 0 -- Striping (!)}
  \begin{columns}
    \begin{column}{0.65\textwidth}
      \textbf{How it works:} data is split into equal-sized blocks (\emph{stripes}) and written
      across all disks in round-robin order.

      \medskip
      Example with 2 disks: Disk\,0 receives blocks A1, A2, A3, A4; Disk\,1 receives A5, A6, A7,
      A8. Both disks are written simultaneously.

      \medskip
      \textbf{Pros:}
      \begin{itemize}
        \item \textbf{Fastest read and write performance} of all RAID levels (full parallelism).
        \item \textbf{100\% storage efficiency} -- no space wasted on redundancy.
      \end{itemize}

      \textbf{Cons:}
      \begin{itemize}
        \item \textbf{Zero fault tolerance}: if even \emph{one} disk fails, all data is lost
              (stripes from the failed disk are missing and cannot be reconstructed).
        \item RAID 0 is technically \emph{not} redundant despite being called RAID.
      \end{itemize}

      \begin{alertblock}{Use case (!)}
        Temporary scratch storage, video editing caches, or any scenario where speed matters and
        data loss is acceptable. Never use for production databases.
      \end{alertblock}
    \end{column}
    \begin{column}{0.3\textwidth}
      \centering
      \textbf{Disk 0} \quad \textbf{Disk 1}\\[4pt]
      A1 \quad A5\\
      A2 \quad A6\\
      A3 \quad A7\\
      A4 \quad A8\\
    \end{column}
  \end{columns}
\end{frame}

\begin{frame}{RAID 1 -- Mirroring (!)}
  \begin{columns}
    \begin{column}{0.65\textwidth}
      \textbf{How it works:} every write operation is duplicated to two (or more) disks
      simultaneously. Disk\,0 and Disk\,1 are always byte-for-byte identical.

      \medskip
      Example with 2 disks: Disk\,0 = A1, A2, A3, A4 and Disk\,1 = A1, A2, A3, A4.

      \medskip
      \textbf{Pros:}
      \begin{itemize}
        \item \textbf{Full redundancy}: if one disk fails, the mirror takes over instantly with no
              data loss and no rebuild time.
        \item \textbf{Fast reads}: the controller can serve read requests from either disk in
              parallel, doubling read throughput.
      \end{itemize}

      \textbf{Cons:}
      \begin{itemize}
        \item \textbf{50\% storage efficiency}: you pay for two disks but only get the capacity of
              one.
        \item Most expensive per usable gigabyte of all RAID levels.
      \end{itemize}

      \begin{examples}{Use case (!)}
        Database servers, operating system drives, any storage where uptime and data availability
        are the top priority.
      \end{examples}
    \end{column}
    \begin{column}{0.3\textwidth}
      \centering
      \textbf{Disk 0} \quad \textbf{Disk 1}\\[4pt]
      A1 \quad A1\\
      A2 \quad A2\\
      A3 \quad A3\\
      A4 \quad A4\\
    \end{column}
  \end{columns}
\end{frame}

\begin{frame}{RAID 5 -- Striping with Distributed Parity (!)}
  \textbf{How it works:} data is striped across $n \geq 3$ disks. For each stripe, a \emph{parity
  block} is computed and stored on one of the disks (the parity location rotates across disks so no
  single disk becomes the bottleneck).

  \medskip
  If one disk fails, its data can be \textbf{mathematically reconstructed} from the remaining
  disks and the parity blocks using XOR operations.

  \medskip
  \textbf{Pros:}
  \begin{itemize}
    \item Good balance of \textbf{performance}, \textbf{storage efficiency}, and
          \textbf{fault tolerance}.
    \item Storage efficiency: $\frac{n-1}{n}$ -- with 4 disks, 75\% of total capacity is usable.
    \item Survives the failure of \textbf{exactly one disk} without data loss.
  \end{itemize}

  \textbf{Cons:}
  \begin{itemize}
    \item \textbf{Slower writes} than RAID 0 or RAID 1 (parity must be computed for every write).
    \item If \textbf{two disks fail simultaneously}, all data is lost (no second parity).
    \item After a disk failure, \textbf{rebuild time} can take hours or days on large arrays --
          during which the remaining disks are under extra stress.
  \end{itemize}

  \begin{block}{Use case (!)}
    The most common RAID level in enterprise NAS and SAN storage arrays where a balance of cost,
    performance, and safety is needed.
  \end{block}
\end{frame}
